\documentclass[../NDCNNAI_main.tex]{subfiles}
\graphicspath{{\subfix{../fig/}}}

% Биологические и искусственные нейроны. Исторический контекст ИНС
% - Основные понятия: биологический нейрон, синапсы, простейшие модели.
% - Переход от биологических нейронных сетей к искусственным.
% - Первая постановка задачи об аппроксимации и обучении сетей.

\begin{document}
\subsection{Основные понятия: биологический нейрон, синапсы, простейшие модели}
Человеческий мозг содержит примерно $8{,}6 \times 10^{10}$ специальных клеток --- нейронов.\\
Каждый нейрон состоит из:
\begin{itemize}
    \item тела клетки (сомы),
    \item множества дендритов (входы),
    \item одного аксона (выход).
\end{itemize}

\begin{figure}[!ht]
    \centering
    \begin{tikzpicture}[scale=0.8]        
        % Дендриты
        \draw (-2.5,1.5) -- (-1,0.6);
        \draw (-2.8,0.5) -- (-1.2,0.2);
        \draw (-2.8,-0.5) -- (-1.2,-0.2);
        \draw (-2.5,-1.5) -- (-1,-0.6);
        \node[left] at (-2.6,0) {Дендриты};
        
        % Сома
        \shade[ball color=gray!40] (0,0) circle (1.2cm);
        \node at (0,0) {Сома};

        % Аксон
        \draw[thick] (1.2,0) -- (4,0);
        \draw (4,0) -- (4.8,0.5);
        \draw (4,0) -- (4.8,0);
        \draw (4,0) -- (4.8,-0.5);
        \node[above] at (2.5,0.2) {Аксон};
        \node[right] at (4.8,0) {Аксонные терминалы};
    \end{tikzpicture}

    \caption{Схематическое изображение биологического нейрона}
\end{figure}

Связь между нейронами осуществляется через \textbf{синапсы}: электрический импульс (потенциал действия $\approx 100$ мВ) вызывает химическую передачу сигнала (время $\approx 1$ мс).
Ключевой принцип --- \textbf{закон <<всё или ничего>>} (Bowditch, 1871): отклик бинарный относительно порога; рефрактерный период ограничивает частоту.

\subsection{Переход от биологических нейронных сетей к искусственным}
Искусственные нейронные сети (ИНС) --- это вычислительные модели, вдохновлённые биологическими нейронными связями и изначально создаваемые для воссоздания человеческого мышления при решении задач.
Сегодня спектр их применения расширился: компьютерное зрение, синтез и обработка речи, переводы на другой язык, обработка медицинских данных, поведение персонажей в играх и многое другое.
Более формально это параметрическое семейство функций, полученных из простых вычислительных единиц (нейронов), сигнал между которыми передаётся с неким весом, и обученное на наборе данных.
Абстракция биологических принципов привела к модели \textbf{перцептрона} (подробное описание далее).

\subsection{Первая постановка задачи об аппроксимации и обучении сетей}
Перцептрон (Rosenblatt, 1957) --- первая модель с правилом обучения: веса корректируются на основе ошибок классификации.
Он реализует линейный классификатор и может различать только линейно разделимые множества.

\begin{definition}[Линейно разделимые множества]
    Множества $A, B \subset \mathbb{R}^n$ \emph{линейно разделимы}, если существуют такие $w \in \mathbb{R}^n$ и $b \in \mathbb{R}$, что
    $$ \langle w, x \rangle + b > 0 \quad \forall x \in A, \quad \langle w, x \rangle + b < 0 \quad \forall x \in B. $$
\end{definition}

\begin{remark}
    Перцептрон не может решить нелинейные задачи (например, представить XOR).
\end{remark}

Для преодоления этого ограничения вводятся многослойные сети с нелинейными функциями активации.
В итоге у нас получается много параметров, которые нужно каким-то образом определить, ставится \textbf{задача обучения:} для набора данных $\{(x_i, y_i)\}_{i=1}^N$ найти веса сети, минимизирующие функцию потерь, например, среднеквадратичную ошибку:
$$MSE(w) = \frac{1}{N} \sum_{i=1}^N \| y(x_i; w) - y_i \|^2$$
Другая интересующая нас проблема --- \textbf{задача аппроксимации:} найти условия, при которых нейронные сети с одним скрытым слоем могут аппроксимировать любую непрерывную функцию на компакте в пространстве $C(K)$ или в $L^p$ пространствах.
\end{document}